% Nouveau EGA V (cote 25): archive cote 25, site page 177 (PDF page 178).
% English translation of the checked French annotated master.
% AI-translated by GPT-6 Astra (OpenAI, ChatGPT Pro mode).
\scanpage{25/177}
\mspage{\hand{4}}

(c) If E has odd rank everywhere, we have $2\overset{\hand{\sim}}{d}(Q)^{\unsure{2}} = d(Q)$.

\var{\xout{Moreover, the ``corrected discriminant'' $\overset{\hand{\sim}}{d}$ thus introduced satisfies}}{}{typewriter strike-out}{ega25_25_177-v0001}

\var{\xout{\shortstack[l]{the\\following conditions}}}{}{typewriter strike-out}{ega25_25_177-v0002} \quad NB It is not reasonable to state the

prope\unsure{r}\var{\over{\illegible{2}}{\unsure{ty}}}{\unsure{ty}}{writing over}{ega25_25_177-v0003} of compatibility with base change without stating

at the same time\hand{,} or even before\hand{,} the property of functoriality under isomorphisms,

given that base change \var{\inserted[typed]{on $\underline{E}$}}{on $\underline{E}$}{insertion}{ega25_25_177-v0004} itself is hardly defined except

up to isomorphism; this would no longer be so if one restricted oneself

to the case $E = \underline{O}_S{}^n$ ,(which would be inconvenient for references).

\var{\xout{\uline{Corollary} 16.7.}}{}{typewriter strike-out}{ega25_25_177-v0006} \quad \uline{Definition} 16.7. The section $\overset{\hand{\sim}}{d}(Q)$ of $\text{dét}(E)^{\hand{\otimes}2}$

will be called the \uline{corrected discriminant} of the quadratic form Q.

\uline{Corollary} 16.8. \unsure{L}\unsure{et} Q,Q' be forms relative to $\underline{E}$, $\underline{E}'$ ; suppose

the parity of the ranks of $\underline{E}$ (resp. $\underline{E}'$) is constant on S. Then

a) If $\underline{E}$ and $\underline{E}'$ do not both have odd rank, we have
\[ \overset{\hand{\sim}}{d}(Q \hand{\oplus} Q') = \overset{\hand{\sim}}{d}(Q)\overset{\hand{\sim}}{d}(Q') \]

b) If $\underline{E}$ and $\underline{E}'$ both have odd rank, we have
\[ \overset{\hand{\sim}}{d}(Q \hand{\oplus} Q') = \hand{2}\,\overset{\hand{\sim}}{d}(Q)\overset{\hand{\sim}}{d}(Q') \]

The verification is trivial.

We can now state the main result of this

number:

\uline{Theorem} 16.9. Let $\underline{E}$ be a loc free Module of finite rank over the prescheme

S, $Q \hand{\in}\hand{\Gamma}\,\uline{Sym}^2(\underline{E})$. The following conditions are equivalent:

(i) Q is smooth.

(ii) The modified d\unsure{s|i}\unsure{s|i}criminant $\overset{\hand{\sim}}{d}(Q) \hand{\in}\hand{\Gamma}\,\text{dét}(\underline{E})^{\hand{\otimes}2}$ is invertible.

(iii) \var{\inserted[typed]{(When $\underline{E}$ has constant rank n)}}{(When $\underline{E}$ has constant rank n)}{insertion}{ega25_25_177-v0007} There is a surjective morphism S' \quad S such that the form $Q_{S'}$

obtained from Q by base change is isomorphic to the standard quadratic

form in n variables.

(i\unsure{i|x}i bis). As in (iii), but assuming S'\hand{$\rightarrow$}S faithfully flat of

finite presentation.

\mspage{\hand{177}}
\tnote{25-177-1}{The two uncertainties in property retain the doubtful r and overwritten ié in propriété; the replacement ty translates the latter ending, with the old two illegible characters retained. The p-to-g retouch in gère has no English spelling effect on hardly. The two parts of Let retain the e/n and n/e doubts in Soient. The canceled s in constants is French adjective agreement and has no English effect. The terms corrected and modified follow the source's respective words. The missing arrow between the two plain S symbols remains unsupplied.}
